% 《大学物理简明教程》第5版第10章解答题10.3.9—10.3.16
% 原书及PDF第273—274页；独立计算、标明单缝次极大近似的有效性。
\subsection{解答题：迈克耳孙干涉仪、单缝与光栅}

\sourceproblem{10.3.9}{273}{273}
\problemtext 迈克耳孙干涉仪一反射镜平移$\Delta d=0.322$ mm，
共观测1024条纹移过视场。求单色光波长。
\workings
反射镜位移导致对应光臂往返光程改变$2\Delta d$。
若移过$N$条纹，则$2\Delta d=N\lambda$，所以
\[
\boxed{\lambda=\frac{2(0.322\times10^{-3})}{1024}
\approx6.29\times10^{-7}\,\mathrm m=629\,\mathrm{nm}}.
\]
\crosscheck 镜子位移和光程差之间有往返因子2；
漏掉会算出约314 nm。

\sourceproblem{10.3.10}{273}{273}
\problemtext 在迈克耳孙干涉仪的一条光臂中插入折射率$n=1.632$的玻璃片，
造成150条纹朝同一方向移过。入射光波长$500$ nm，求片厚。
\workings
插入片后单程光程多$(n-1)t$，
仪器光线往返两次，总光程差改变
$2(n-1)t$；故
\[
\boxed{t=\frac{N\lambda}{2(n-1)}
=\frac{150(500\times10^{-9})}{2(1.632-1)}
\approx5.93\times10^{-5}\,\mathrm m
=59.3\,\mu\mathrm m}.
\]
\crosscheck 增厚、增大折射率或增加移动条数三者正相关，
与公式一致。

\sourceproblem{10.3.11}{273}{273}
\problemtext 对垂直入射的多缝光栅，透光缝宽$a$，
不透光区宽$b$，讨论三种光栅常量比例
$(1)\ a+b=2a$、$(2)\ a+b=3a$、
$(3)\ a+b=4a$时哪些主极大级数会因单缝包络而缺级。
\workings
光栅主极大$d\sin\theta=k\lambda$，其中$d=a+b$；
单缝零点$a\sin\theta=m\lambda$。
两式同时成立须
\[
k=\frac{d}{a}\,m,\quad m=1,2,\ldots.
\]
因此
\[
\boxed{\begin{array}{c|c}
d/a&\text{消失的非零主极大级数}\\ \hline
2& k=\pm2,\pm4,\pm6,\ldots\\
3& k=\pm3,\pm6,\pm9,\ldots\\
4& k=\pm4,\pm8,\pm12,\ldots
\end{array}}
\]
\crosscheck $k=0$中心主极大不会因为单缝零点而缺失；
是否能观察高阶级数还须满足$|\sin\theta|<1$。

\sourceproblem{10.3.12}{273}{273}
\problemtext 用未知波长$\lambda$单缝衍射时第3级暗纹位置，
与波长$600$ nm光的第2条次极大亮纹位置恰好重合，求未知波长。
\workings
第3级暗纹的光程差满足$a\sin\theta=3\lambda$。
常用\emph{半波带近似}把单缝第2侧明纹放在
$a\sin\theta\simeq(2+\tfrac12)\lambda_0$，
$\lambda_0=600$ nm，因此
\[
\boxed{\lambda_{\rm approx}
\simeq\frac{2.5(600)}{3}\,\mathrm{nm}=500\,\mathrm{nm}}.
\]
更精细的夫琅禾费次极大满足
$\tan\beta=\beta$，
第2个非中心正侧极大根为$\beta\simeq7.725$，
对应$a\sin\theta\simeq2.459\lambda_0$，于是
\[
\lambda_{\rm exact}\simeq2.459(600)/3\approx492\,\mathrm{nm}.
\]
\crosscheck \textbf{模型精度必须注明}：
标准教材以“奇数个半波带”近似求次明纹会给500nm，
但该位置不是严格的$\tan\beta=\beta$极值。
不能把500nm标成精确值；根的级次应以零点之间的
\emph{第2条侧极大}为准。

\sourceproblem{10.3.13}{273}{273}
\problemtext 单缝宽$a=0.10$ mm，原装置凸透镜焦距$f=50$ cm，
波长$500$ nm的绿光垂直入射。将整个装置浸入折射率$n=1.33$的水中，
求中央衍射明纹的\emph{半角宽度}。
\workings
题问的是角度而不是焦平面的线宽。
水中光波频率不变、波长
$\lambda'=\lambda_0/n$，
单缝第一暗纹角满足
$a\sin\theta_1=\lambda'$。
因此中央明纹从中心到第一暗纹的半角宽度为
\[
\boxed{\theta_{\frac12}
=\arcsin\frac{\lambda_0}{na}
=\arcsin\frac{500\times10^{-9}}{1.33(0.10\times10^{-3})}
\approx3.76\times10^{-3}\,\mathrm{rad}}
\]
即约$\boxed{0.215^\circ}$。
\crosscheck 焦距$f=50$ cm对\emph{角宽度}不起作用；
若另外换算假定$f$固定时焦面线半宽才有
$f\theta\simeq1.88$ mm。不能将mm作为原题半角宽度的单位。

\sourceproblem{10.3.14}{273--274}{273--274}
\problemtext 橙黄色平行光垂直入射宽$a=0.60$ mm的单缝，
后方凸透镜焦距$f=40.0$ cm；
屏上距中心$1.40$ mm处的$P$点为侧明纹。
求大致波长、明纹级次及从$P$点观察的半波带数。
\workings
旁轴近似给该方向单缝光程差
\[
\delta=a\sin\theta\simeq a\frac yf
=(0.60\times10^{-3})\frac{1.40\times10^{-3}}{0.40}
=2.10\times10^{-6}\,\mathrm m.
\]
橙黄色波长通常约$600$ nm，
故$\delta/\lambda\approx3.5$，
按照半波带近似，$P$对应\emph{第3个侧次极大}，
且光程差约$7\lambda/2$，由此
\[
\boxed{\lambda_{\rm approx}\simeq
\frac{\delta}{7/2}\approx600\,\mathrm{nm}},\quad
\boxed{\text{侧明纹级次约为第3级}},\quad
\boxed{\text{约7个半波带}}.
\]
\emph{更精细地}，单缝第3侧极大的根
$\beta\simeq10.904$，
光程差系数$\beta/\pi\simeq3.4709$，
由同一屏位置可反推$\lambda\approx605$ nm。
\crosscheck 原文只称“橙黄色”，未独立给出精确波长；
单靠“$P$点是一条明纹”会有多个可能级次。
上述600nm的选择额外使用\emph{颜色范围及第3侧明纹}
的近似判别；不应声称题目在数学上无歧义唯一确定波长。

\sourceproblem{10.3.15}{274}{274}
\problemtext 波长$590$ nm钠黄光垂直入射每毫米有500条刻线的光栅，
求能观察到的最高主极大级数。
\workings
光栅常量
$d=1/500\,\mathrm{mm}=2.00\,\mu\mathrm m$。
主极大满足$d\sin\theta=k\lambda$，
且可观测传播方向要求$|\sin\theta|\le1$。
\[
\frac d\lambda=\frac{2000}{590}\approx3.39,
\qquad\boxed{k_{\max}=3}.
\]
\crosscheck 题未给单缝宽度，因此不额外假定第3级缺级。

\sourceproblem{10.3.16}{274}{274}
\problemtext 两波长$\lambda_1=600$ nm、$\lambda_2=400$ nm的平行单色光
同时垂直入射光栅。屏上距零级明纹$y=5$ cm处，
前者第$k$级极大与后者第$k+1$级重合，
透镜焦距$f=50$ cm，求$k$和光栅常量$d$。
\workings
共同衍射角要求
$k\lambda_1=(k+1)\lambda_2$，
故
\[
\boxed{k=\frac{\lambda_2}{\lambda_1-\lambda_2}
=\frac{400}{600-400}=2}.
\]
用焦平面几何$\tan\theta=y/f=0.10$，
则$\sin\theta=y/\sqrt{f^2+y^2}\approx0.09950$。
\[
\boxed{d=\frac{k\lambda_1}{\sin\theta}
=\frac{2(600\,\mathrm{nm})}{0.09950}
\approx12.1\,\mu\mathrm m}.
\]
若采用小角近似$\sin\theta\simeq y/f$，得$d\simeq12.0\,\mu$m。
\crosscheck 用的是光栅\emph{常量}而不是透光缝宽，
且两束光的级数相差1而非相同。
